\chapter*{How to Read This Book}
\addcontentsline{toc}{chapter}{How to Read This Book}
\markboth{How to Read This Book}{How to Read This Book}

This book teaches you to design digital devices. It is written for people
who are just starting out in digital electronics: school-level mathematics
and a wish to understand how computers, drone controllers, and any other
digital devices ``think'' are enough.

\section*{The Idea of the Book: Design, Run, Observe}

You cannot learn a digital circuit just by reading about it. That is why in
every chapter you repeat the same cycle:

\begin{enumerate}
  \item \textbf{Design.} Build a circuit in the Logra Designer or describe it
    in Code.
  \item \textbf{Run.} Set the inputs and start the simulation.
  \item \textbf{Observe.} Compare the waveform with the table in the book,
    look at the schematic on the RTL tab, and explain why the result came out
    exactly this way.
\end{enumerate}

Logra is the environment you work in from the first page. It consists of the
tabs Learn, Code, Designer, Pins, RTL, Waveforms, Lab, Gowin, Chip, and
Connect. In the sidebar they are collected into the groups ``Learn'',
``Create'', ``Check'', and ``Real board''. You can start working with it right
away: the first chapters rely only on the Designer, Waveforms, and RTL.

\section*{How a Chapter Is Organized}

Each chapter is a small ``sprint''. It begins with a bridge: what we can
already do and what question is still open. Then comes a ``Build it first''
box with a task in Logra that is worth doing before you read the theory: this
way you meet the theory already holding a question. Every section ends with a
\textbf{``Do it yourself''} task: what to build in Logra, what result to
expect, and how to check it. At the end of the chapter come common mistakes,
questions and problems, and a bridge to the next topic.

\begin{infobox}[The Chapter Rule]
Do not move on to the next section until the waveform agrees with the table.
If the signals on the Waveforms tab disagree with the truth table, then the
mistake is either in the circuit or in the table. Find it now: later it will
only pile up.
\end{infobox}

\subsection*{Difficulty Levels}

Tasks and problems are marked with levels:
\stars{1}~--- \emph{repeat}: do it by the example and check the result;
\stars{2}~--- \emph{modify}: change a ready-made circuit or find a mistake in it;
\stars{3}~--- \emph{project}: solve the problem on your own and accept the
result against a given criterion. If it is hard, go back one level; that is
normal.

\section*{The Running Project ``Security Panel''}

Throughout the book you build a single device: a security panel with door and window
sensors, a night mode, an arming key, a code lock, a siren, and an indicator.
Each chapter adds a new part to it (a sensor as a bit, the alarm formula,
sensor voting, zone selection, a noise threshold, the lock code, alarm
memory, a time counter, the arming state machine), and the panel always
stays working. By the end of the theory line you have a finished device, and
in the chapters on the board and the projects it ends up on an FPGA. The stages are collected in
Appendix~\ref{app:logramap}.

\section*{How the Book Is Organized}

The book is a single continuous line, with no division into parts. Theory,
Verilog, and work in Logra go together: at every step you design, run, and
observe.

\begin{description}
  \item[Start.] What an FPGA is and how it differs from a processor; the
    first launch of Logra: the tabs and the ``design, run, observe'' cycle.
  \item[Sprints by topic.] Signals; number systems and codes; logic gates and
    Boolean algebra; the Verilog language; combinational circuits;
    minimization (Karnaugh maps, the Quine---McCluskey method); decoders and
    encoders; multiplexers; arithmetic blocks; flip-flops; registers and
    counters; synchronous circuits and finite-state machines; state machine
    design and FIFO queues. In every chapter the theory, the Verilog code, and
    the circuit in Logra are studied together.
  \item[Board, flow, physics, analog.] What is inside an FPGA (LUTs,
    flip-flops, routing); the Sipeed Tang Nano 20K board; the path from text
    to a bitstream; the electrical basics of CMOS and buses; DAC, ADC, and
    PWM; the bridge from a circuit to Verilog.
  \item[Projects and systems.] Seven complete projects; simulation and
    testbenches; clocking, timing analysis, and clock domain crossing; the
    UART, SPI, and I$^2$C interfaces; SDRAM memory; your own processor that
    runs a program; assignments for independent work.
\end{description}

\begin{figure}[H]
\centering
\begin{tikzpicture}[node distance=5mm and 4mm, t/.style={block, minimum width=28mm, minimum height=15mm, font=\sffamily\footnotesize}]
  \node[t] (a) {Start:\\FPGA and Logra};
  \node[t, right=of a] (b) {Signals,\\gates};
  \node[t, right=of b] (c) {Verilog,\\combinational\\blocks};
  \node[t, right=of c] (d) {Flip-flops,\\registers};
  \node[t, right=of d] (e) {State\\machines};
  \node[hblock, minimum width=28mm, minimum height=15mm, font=\sffamily\footnotesize, below=9mm of e] (f) {Board,\\flow};
  \node[hblock, minimum width=28mm, minimum height=15mm, font=\sffamily\footnotesize, left=of f] (g) {Physics,\\analog};
  \node[hblock, minimum width=28mm, minimum height=15mm, font=\sffamily\footnotesize, left=of g] (h) {Bridge to\\Verilog};
  \node[oblock, minimum width=28mm, minimum height=15mm, font=\sffamily\footnotesize, left=of h] (i) {Projects,\\UART, clocks};
  \node[oblock, minimum width=28mm, minimum height=15mm, font=\sffamily\footnotesize, left=of i] (j) {SDRAM, CPU,\\assignments};
  \foreach \u/\v in {a/b, b/c, c/d, d/e, e/f, f/g, g/h, h/i, i/j} \draw[arr] (\u) -- (\v);
\end{tikzpicture}
\caption{The path we will travel}
\end{figure}

If you already know logic, flip-flops, and state machines, you can go
straight to the chapters on the board and the projects (starting with
Chapter~\ref{ch:inside}) and return to the sprint chapters as a reference:
they are cross-referenced. The correspondence between chapters and Logra
lessons and examples is collected in Appendix~\ref{app:logramap}.

All the source texts of the examples are stored next to the book, in the
folder \nolinkurl{textbook/code}. You can open them, copy them, and load them
onto the board right away.

\begin{infobox}[Conventions]
\begin{itemize}
  \item Blue ``Reference'' boxes contain useful details.
  \item Orange ``Warning!'' boxes warn about common mistakes.
  \item Purple ``Key idea'' boxes sum up a section.
  \item Green ``Task'' and ``Do it yourself'' boxes suggest something to do
    on your own in Logra.
  \item Light-blue ``Build it first'' boxes give a task before the theory.
  \item Italic ``Looking ahead'' notes mark rare glimpses into later chapters.
\end{itemize}
\end{infobox}
